\documentclass[11pt,a4paper]{article}

\usepackage{amsmath,amssymb,amsthm}
\usepackage{graphicx}
\usepackage{booktabs}
\usepackage{hyperref}
\usepackage[margin=1in]{geometry}
\usepackage{algorithm}
\usepackage[noend]{algpseudocode}

\newtheorem{definition}{Definition}
\newtheorem{theorem}{Theorem}
\newtheorem{proposition}{Proposition}

\title{Symbolic Mechanism Discovery via Best-First Search\\in Mathematical Operation Space}

\author{Andrew H. Bond\\
\small Department of Computer Engineering, San Jose State University\\
\small Senior Member, IEEE \quad ORCID: 0009-0003-2599-6158\\
\small \texttt{andrew.bond@sjsu.edu}}

\date{\today}

\begin{document}
\maketitle

\begin{abstract}
We introduce a framework for automated discovery of interpretable
predictive formulas from labeled scientific data using best-first
search through a space of mathematical operations. Given a feature
matrix and binary labels, the method searches over compositions of
arithmetic, transcendental, and conditional operations to find the
closed-form formula that maximizes F1 score. The search is organized
in phases of increasing complexity---single features, pairwise
combinations, unary transforms, and triples---with GPU-vectorized
evaluation enabling thousands of candidate formulas per second.

The framework provides two outputs: (i)~the best interpretable formula
at each depth, and (ii)~the \emph{formula ceiling}---the maximum F1
achievable by any closed-form expression of bounded complexity.
Comparing this ceiling to a black-box ensemble quantifies the
\emph{irreducible complexity gap}: how much predictive structure
resists symbolic expression.

Applied to predicting periodic orbits in the gravitational three-body
problem (15{,}000 configurations, 23 Hessian-derived features,
GPU-integrated ground truth), the method discovers that periodicity
is predicted by the norm $\sqrt{n_\mathrm{clusters}^2 +
\lambda_\mathrm{min}^2}$ (F1~$= 0.48$). The ceiling is robust across
all tested operations. A gradient-boosted ensemble reaches
F1~$= 0.86$, yielding an irreducible complexity gap of 0.38---meaning
38\% of the predictive structure in three-body dynamics cannot be
expressed as a depth-4 formula.

Code and data: \url{https://github.com/ahb-sjsu/astar-symbolic-discovery}.
\end{abstract}

\section{Introduction}

A central goal of computational science is to discover interpretable
laws from data. This problem has been addressed by genetic symbolic
regression~\cite{schmidt2009distilling, cranmer2023interpretable},
neural equation discovery~\cite{udrescu2020ai}, and classical AI
systems~\cite{langley1987scientific}. Genetic programming explores
large formula spaces but provides no optimality guarantee at a given
depth and often produces bloated expressions. Neural methods discover
accurate but opaque representations. Classical systems require
domain-specific search heuristics.

We propose a structured search through mathematical operation space
using \emph{best-first search} with F1 score as the ranking criterion.
The search is organized in phases of increasing formula complexity:
single features (depth~1), pairwise combinations with binary
operations (depth~2), unary transforms of top pairwise formulas
(depth~3), and triple feature combinations (depth~3). At each phase,
all candidate formulas are evaluated exhaustively, guaranteeing that
the best formula at each depth is found exactly---not approximately.

This phased exhaustive approach avoids the admissibility issues of
heuristic A* search on formula spaces (where defining an admissible
heuristic for future F1 improvement is nontrivial) while preserving
the key practical benefit: the best formula at each depth is
identified with certainty.

Our framework provides two outputs that existing methods do not:
\begin{enumerate}
\item The \textbf{best interpretable formula} at each complexity level,
found by exhaustive search within that level.
\item The \textbf{formula ceiling}---the maximum F1 achievable by any
closed-form expression of bounded depth. Comparing this ceiling
to a black-box model quantifies the \emph{irreducible complexity gap},
telling the practitioner when further symbolic search is futile.
\end{enumerate}

\section{Method}

\subsection{Problem Formulation}

\begin{definition}[Symbolic Prediction Problem]
Given a feature matrix $X \in \mathbb{R}^{N \times d}$ and binary
labels $y \in \{0, 1\}^N$, find a formula
$f: \mathbb{R}^d \to \mathbb{R}$ and threshold $\tau \in \mathbb{R}$
that maximize $F_1(\hat{y}, y)$ where
$\hat{y}_i = \mathbf{1}[f(X_i) > \tau]$, subject to
$\mathrm{depth}(f) \leq D$.
\end{definition}

\subsection{Operation Registry}

The search space is defined by three operation classes:

\textbf{Leaf nodes} (depth 0): Each of the $d$ features $x_j$.

\textbf{Unary operations} (depth $+1$): $\{\log|\cdot|, \sqrt{|\cdot|},
(\cdot)^2, 1/(\cdot), -(\cdot), |\cdot|, \sigma(\cdot),
\tanh(\cdot)\}$ (8 operations).

\textbf{Binary operations} (depth $+1$): $\{+, -, \times, \div,
\max, \min, \mathrm{hypot}, (\cdot - \cdot)^2,
\mathrm{harmonic\text{-}mean}, \mathrm{geometric\text{-}mean}\}$
(10 operations).

\subsection{Phased Exhaustive Search}

Rather than a single tree search, we organize the exploration into
four phases (Algorithm~\ref{alg:search}). Each phase exhaustively
evaluates all formulas at its complexity level.

\begin{algorithm}[t]
\caption{Phased Exhaustive Symbolic Search}
\label{alg:search}
\begin{algorithmic}[1]
\Require Features $X \in \mathbb{R}^{N \times d}$, labels $y$,
operations $\mathcal{U}$ (unary), $\mathcal{B}$ (binary)
\Ensure Best formula and F1 at each depth
\Statex
\Statex \textbf{Phase 1:} Single features (depth 1)
\For{$j = 1, \ldots, d$}
    \State Evaluate $F_1^*(x_j)$ via threshold sweep
\EndFor
\State $R_1 \gets$ top-$K$ features by F1
\Statex
\Statex \textbf{Phase 2:} Pairwise combinations (depth 2)
\For{$x_i, x_j \in \{1,\ldots,d\}$, $b \in \mathcal{B}$}
    \State Evaluate $F_1^*(b(x_i, x_j))$ via threshold sweep
\EndFor
\State $R_2 \gets$ top-$K$ pairwise formulas
\Statex
\Statex \textbf{Phase 3:} Unary transforms (depth 3)
\For{$f \in R_2$, $u \in \mathcal{U}$}
    \State Evaluate $F_1^*(u(f))$ via threshold sweep
\EndFor
\State $R_3 \gets$ top-$K$ depth-3 formulas
\Statex
\Statex \textbf{Phase 4:} Triple combinations (depth 3)
\For{$x_i, x_j, x_k \in R_1$, pattern $p$}
    \State Evaluate $F_1^*(p(x_i, x_j, x_k))$ via threshold sweep
\EndFor
\State \Return best formula across all phases
\end{algorithmic}
\end{algorithm}

\textbf{Threshold sweep.} For each candidate formula $f$, we compute
$v_i = f(X_i)$ for all $N$ samples, then evaluate F1 at percentile
thresholds $\{20, 30, 40, \ldots, 95\}$ of the value distribution.
Both $v > \tau$ and $v < \tau$ predictors are tested. The best
threshold is retained. This sweep is vectorized and adds negligible
cost.

\textbf{GPU acceleration.} Features are stored as GPU arrays (CuPy).
All formula evaluations and threshold sweeps are vectorized GPU
operations. Phase~2 evaluates $d^2 \times |\mathcal{B}|$ formulas
(e.g., $23^2 \times 10 = 5{,}290$) in seconds.

\textbf{Completeness within each phase.} Phases~1--2 are exhaustive
over all feature (pairs) and operations. Phase~3 applies all unary
transforms to the top-$K$ results from Phase~2. Phase~4 tests
multiple combination patterns over the top-$K$ single features.
The search is complete at each phase for the specified formula
structure; no formulas of that structure are missed.

\subsection{Scalability}

The formula count per phase is:
Phase~1: $O(d)$;
Phase~2: $O(d^2 \cdot |\mathcal{B}|)$;
Phase~3: $O(K \cdot |\mathcal{U}|)$;
Phase~4: $O(\binom{K}{3} \cdot P)$ where $P$ is the number of
triple-combination patterns.

For $d = 23$, $|\mathcal{B}| = 10$, $|\mathcal{U}| = 8$, $K = 50$,
$P = 9$, the total is approximately 6{,}200 formulas---evaluated in
under one minute on a single GPU. Depth~5--6 formulas (Phase~2
composed with Phase~2) would require $O(d^4 \cdot |\mathcal{B}|^2)
\approx 28$ million formulas, which is tractable with GPU
vectorization but would take hours. We do not explore beyond depth~4
in this work.

\subsection{Formula Ceiling and Irreducible Complexity Gap}

\begin{definition}[Formula Ceiling]
The formula ceiling $C_D$ at depth $D$ is the maximum F1 achieved
by any formula of depth $\leq D$ found across all phases:
$C_D = \max_{\text{phases}} F_1^*$.
\end{definition}

\begin{definition}[Irreducible Complexity Gap]
Given a black-box model with F1 score $F_1^{\mathrm{bb}}$ trained on
the same features, the irreducible complexity gap at depth $D$ is
$G_D = F_1^{\mathrm{bb}} - C_D$.
\end{definition}

$G_D > 0$ means the data contains predictive structure that cannot
be captured by depth-$D$ formulas. $G_D$ is a property of the data
and the feature set, not of the search algorithm. If $C_D$ converges
(does not increase with additional operations or deeper phases),
further symbolic search is futile.

\begin{proposition}
$C_D$ is non-decreasing in $D$ and bounded above by
$F_1^{\mathrm{Bayes}}$. If $C_{D+1} = C_D$ for a sufficiently
rich operation set, the ceiling has converged.
\end{proposition}

\section{Case Study: Three-Body Periodic Orbit Prediction}

\subsection{Setup}

We compute 23 features from the Hessian of the gravitational
three-body Hamiltonian at 15{,}000 random configurations of
equal-mass bodies ($m_1 = m_2 = m_3 = 1$) in Jacobi coordinates.
Features include: numerical matrix rank (threshold $\epsilon = 10^{-6}$),
potential block eigenvalues and their ratios, frequency cluster counts,
spectral entropy, Hamiltonian energy, body separation ratios, and
derived quantities. Ground truth is obtained by symplectic leapfrog
integration ($\Delta t = 0.005$, $8 \times 10^4$ steps) on an
NVIDIA GV100 GPU, with periodicity defined as phase-space return
to within 5\% of initial extent. Batch size is maximized
automatically using the batch-probe library~\cite{batchprobe}.

\subsection{Search Results}

Table~\ref{tab:phases} summarizes the results across all four phases.

\begin{table}[h]
\caption{Formula search results by phase (15{,}000 configs, 23 features).}
\label{tab:phases}
\begin{tabular}{llrl}
\toprule
Phase & Best formula & F1 & Formulas tested \\
\midrule
1 (single) & \texttt{n\_clusters} & 0.35 & 23 \\
2 (pairwise) & \texttt{n\_clust + $\lambda_\mathrm{min}$} & 0.48 & 5{,}060 \\
3 (unary) & \texttt{hypot(n\_clust, $\lambda_\mathrm{min}$)} & 0.48 & 400 \\
4 (triple) & (none above 0.30) & --- & 504 \\
\midrule
\textbf{Ceiling} & & \textbf{0.48} & \textbf{5{,}987} \\
\bottomrule
\end{tabular}
\end{table}

The ceiling converges at Phase~2: unary transforms (Phase~3) and
triple combinations (Phase~4) provide no improvement. All 20 top
formulas are algebraically equivalent to
$\alpha \cdot n_\mathrm{clusters} + \beta \cdot \lambda_\mathrm{min}$
for various $\alpha, \beta$. The Euclidean norm variant
\begin{equation}
\hat{y} = \mathbf{1}\!\left[
\sqrt{n_\mathrm{clusters}^2 + \lambda_\mathrm{min}^2} > \tau
\right]
\end{equation}
is representative of the equivalence class.

\subsection{Irreducible Complexity Gap}

\begin{table}[h]
\caption{Formula ceiling vs.\ black-box models (51{,}000 configs, 3 mass ratios).}
\label{tab:gap}
\begin{tabular}{lcc}
\toprule
Model & F1 & Gap from ceiling \\
\midrule
Random baseline (5\% rate) & 0.07 & --- \\
Separation ratio $r_2/r_1 > 10$ & 0.29 & --- \\
Best depth-4 formula (Eq.~1) & 0.48 & 0.00 (= ceiling) \\
Decision tree (depth 5) & 0.58 & $-$0.10 \\
Gradient boosting (100 trees) & 0.86 & $-$0.38 \\
\bottomrule
\end{tabular}
\end{table}

The irreducible complexity gap is $G_4 = 0.86 - 0.48 = 0.38$.
The decision tree (F1~$= 0.58$) captures some of the gap via
piecewise-constant splits; the remaining 0.28 requires the nonlinear
interactions of an ensemble. This means that 38\% of the predictive
structure in three-body periodicity resists expression as any
depth-4 formula over Hessian features.

\subsection{Physical Interpretation}

The discovered formula combines two independent necessary conditions:
\begin{itemize}
\item $n_\mathrm{clusters} \geq 2$: the eigenvalue frequency spectrum
separates into fast (inner binary) and slow (outer body) groups,
indicating timescale separation.
\item $\lambda_\mathrm{min} > 0$: the smallest Hessian eigenvalue is
nonzero, indicating the system is gravitationally bound (not
escaping).
\end{itemize}
The Euclidean norm indicates these contribute independently---there
is no synergy or interaction between them. Periodicity requires
\emph{both} separated timescales \emph{and} sufficient binding.

\subsection{Comparison to PySR}

To contextualize our method, we note how it differs from PySR/genetic
symbolic regression~\cite{cranmer2023interpretable} on this task:

\textbf{Completeness vs.\ exploration.} Our phased search evaluates
all $d^2 \times |\mathcal{B}|$ pairwise formulas exhaustively.
Genetic methods sample this space stochastically and may miss the
global optimum at a given depth, but can explore deeper formula trees
(depth 10+) that our exhaustive approach cannot reach.

\textbf{Formula ceiling.} Our framework explicitly computes the
ceiling and the gap to a black-box model. PySR produces formulas of
increasing complexity but does not systematically measure when
further complexity stops helping.

\textbf{Speed.} Our GPU-vectorized evaluation of $\sim$6{,}000
formulas completes in under one minute. PySR typically requires
minutes to hours depending on population size and generations. For
depth $\leq 4$ with moderate feature counts ($d < 30$), exhaustive
search is faster; for deeper or wider searches, genetic methods
are more practical.

A direct head-to-head comparison on standard symbolic regression
benchmarks (e.g., AI Feynman~\cite{udrescu2020ai}) is left for
future work; our current results demonstrate the method on a novel
scientific prediction task rather than on benchmark recovery.

\section{Discussion}

\subsection{When to Use This Method}

The phased exhaustive approach is most valuable when:
(i)~the feature count is moderate ($d < 50$, so $d^2$ is tractable);
(ii)~the practitioner wants a \emph{guaranteed} best formula at each
depth, not a stochastic approximation;
(iii)~the formula ceiling and irreducible complexity gap are of
interest (e.g., to decide whether a problem admits a simple law or
requires a complex model).

For large feature spaces ($d > 100$) or deep formula trees
(depth $> 5$), genetic methods or neural symbolic regression are
more appropriate.

\subsection{The Ceiling as a Scientific Diagnostic}

The formula ceiling is not just a performance metric---it is a
statement about the data. A converged ceiling means the relationship
between features and labels has a \emph{fundamental} nonlinearity
that no short formula can capture. In our three-body case study,
the ceiling of 0.48 at depth~4 with the gap of 0.38 to gradient
boosting tells us that three-body periodic orbit prediction is
genuinely complex: it depends on nonlinear feature interactions that
the underlying physics does not reduce to a compact law at the level
of Hessian features.

This diagnostic is novel. Existing symbolic regression methods report
the best formula found but do not quantify \emph{how much better}
a more complex model could do. The ceiling provides a principled
stopping criterion: if $C_{D+1} = C_D$ and $G_D > 0$, stop searching
for formulas and use an ensemble.

\subsection{Limitations}

The search is exhaustive only within each phase's structure.
Phase~2 (pairwise) is complete; Phase~3 applies unary transforms
only to top-$K$ pairwise results, not all; Phase~4 tests only
pre-specified triple patterns. A formula with structure not matching
any phase template (e.g., $\log(x_1 + x_2) \cdot x_3 / x_4$) will
be missed.

The method does not discover features---it evaluates pre-computed
ones. The quality of the discovered formula depends entirely on the
feature engineering.

The F1 metric weights precision and recall equally. Domain-specific
objectives (e.g., prioritizing recall for safety applications) could
change the optimal formula.

\section{Conclusion}

We have introduced a phased exhaustive search framework for symbolic
mechanism discovery that provides two novel contributions:
guaranteed-best formulas at each complexity level, and the formula
ceiling as a diagnostic for irreducible complexity. Applied to the
three-body problem, the method discovers that periodicity depends on
two Hessian features (frequency cluster count and minimum eigenvalue)
with a hard ceiling at F1~$\approx 0.48$, while ensemble methods
reach 0.86. The 38\% irreducible complexity gap quantifies the
fundamental nonlinearity of three-body dynamics at the level of
Hessian features.

The method is domain-general, requiring only a feature matrix and
binary labels. It is practical for feature spaces up to $d \approx 50$
and formula depths up to 4, with GPU acceleration enabling evaluation
of $\sim$6{,}000 formulas in under one minute.

\begin{thebibliography}{9}

\bibitem{schmidt2009distilling}
M.~Schmidt and H.~Lipson, ``Distilling free-form natural laws from
experimental data,'' \emph{Science}, vol.~324, pp.~81--85, 2009.

\bibitem{udrescu2020ai}
S.-M.~Udrescu and M.~Tegmark, ``AI Feynman: A physics-inspired
method for symbolic regression,'' \emph{Science Advances}, vol.~6,
no.~16, eaay2631, 2020.

\bibitem{langley1987scientific}
P.~Langley, H.~A.~Simon, G.~L.~Bradshaw, and J.~M.~Zytkow,
\emph{Scientific Discovery: Computational Explorations of the
Creative Processes}. MIT Press, 1987.

\bibitem{cranmer2023interpretable}
M.~Cranmer, ``Interpretable machine learning for science with PySR
and SymbolicRegression.jl,'' \emph{arXiv:2305.01582}, 2023.

\bibitem{batchprobe}
A.~H.~Bond, batch-probe: GPU memory batch-size probing, v0.3.0
(2026). \url{https://pypi.org/project/batch-probe/}

\end{thebibliography}

\end{document}
